\section{Discussion}\label{sec:discussion}
The main structural result is a separation of physical roles. The timelike scalar sector provides an effective cold cosmological component, the noncanonical spatial scalar function provides the tested quasistatic MOND response, the nonseparable current-linked operator fixes the radiation-era scalar damping, and the auxiliary Type-II constraint sector supplies the reconstructed late-time acceleration. The new operator is therefore not an extra phenomenological function used to chase the cosmological likelihood; it is fixed by the conserved scalar current and accepted on stability grounds before the final cosmological optimization.

The statistical comparison should be interpreted conservatively. A lower $\chi^2$ and lower AIC show that the zero-particle-CDM realization is competitive with the matched $\LCDM$ reference within the chosen likelihood stack, while the small positive $\Delta$BIC reflects the larger parameter penalty in the adopted effective-$N$ convention. Neither statistic is a Bayes factor. Likewise, the particle-CDM profile shows that zero particle CDM is not disfavored at the $\Delta\chi^2=1$ level in the sampled profile, but the profile minimum is not exactly at zero. The data therefore support a non-requirement statement, not a detection of the absence of particle dark matter.

The effective-field-theory result is similarly scoped. Canonical normalization places the lowest tested cubic derivative scale many orders of magnitude above the cosmological frequencies and momenta used in the calculation, and the first-order scalar has a regular Dirac reduction. This establishes perturbative control over the tested classical domain. It is not a proof of a fundamental ultraviolet completion.

Several limitations remain relevant for future empirical tests. First, the cosmological analysis is linear; nonlinear large-scale structure, weak lensing, cluster mergers, and void statistics require a dedicated nonlinear solver. Second, the preferred-frame expressions reported in Appendix~\ref{app:dirac} are a screened limiting specialization; a complete moving-source finite-field 1PN derivation remains outside the present analysis. Third, the galactic calculation demonstrates the static no-halo limit on the retained SPARC sample, but cluster-scale morphology and environmental effects require separate analysis. These limitations are intentionally not hidden behind the successful linear-cosmology fit.

The combined construction also differs conceptually from adding a dark-energy fluid to AeST. The VCDM acceleration variable is auxiliary and does not supply another local propagating scalar mode \cite{DeFeliceMaedaMukohyamaPookkillath2022}. Conversely, the effective cold cosmological component is carried by the scalar--vector sector rather than by particle CDM. The final model thus separates the two conventional dark-sector roles into constrained gravitational dynamics rather than replacing both with a single invisible material component.
